\documentclass[journal]{IEEEtran}
\usepackage[spanish]{babel}
\usepackage[utf8]{inputenc}
\usepackage{enumerate}
\usepackage{graphicx}
\usepackage{booktabs}
\usepackage{physics}

\hyphenation{op-tical net-works semi-conduc-tor}


\begin{document}

### Bitácora de un Físico: Exploraciones en Mecánica Vectorial

#### Incertidumbre por Redondeo: Cifras Significativas

---

**Introducción a la Incertidumbre y las Cifras Significativas**

Hoy he comenzado una serie de experimentos enfocados en comprender la incertidumbre por redondeo y las cifras significativas en las mediciones físicas. El redondeo, como descubrí, es una fuente inevitable de error en cualquier medición. Este entendimiento me llevó a considerar cómo una medida hecha con un instrumento de medición puede ser vista como el resultado de redondear el valor de otra medición más precisa. La relación entre la incertidumbre absoluta y relativa en la medición de longitudes y áreas se volvió una pieza clave en mi investigación.

**Medición y Redondeo: Raíz Cuadrada de 3**

Para comprender mejor la incertidumbre en las cifras significativas, comencé con la raíz cuadrada de 3. Redondeé este valor con distintas precisiones:
- 1.7 ± 0.05
- 1.73 ± 0.005
- 1.732 ± 0.0005

Cada nuevo decimal añadido reducía la incertidumbre, haciéndome consciente de la importancia de las cifras significativas en una medición precisa.

**Medición y Redondeo: Raíz Cuadrada de 12 y 15**

Seguí el mismo proceso con la raíz cuadrada de 12:
- 3.46410 ± 0.00005

Y luego con la raíz cuadrada de 15:
- 3.9 ± 0.05
- 3.87 ± 0.005
- 3.872983346 ± 0.0000000005

Estas mediciones confirmaron que la precisión mejora con más cifras significativas, y la incertidumbre disminuye considerablemente.

**Gráficas y Ajustes Lineales**

Graficar las parejas de valores (x, y) me permitió observar las relaciones lineales y ajustar las ecuaciones de las rectas. Para los valores (1, 8), (2, 13), (3, 18), (4, 23), (5, 28), obtuve la ecuación y = 5x + 3. Este ejercicio me mostró cómo las incertidumbres pueden influir en la interpretación de los datos graficados.

**Medición del Diámetro de una Moneda**

Usé diferentes instrumentos para medir el diámetro de una moneda:
- Regla: 2.5 cm
- Vernier: 2.54 cm
- Tornillo micrométrico: 2.545 cm

Cada método ofrecía distintos niveles de precisión y me permitió ver de primera mano cómo las cifras significativas afectan la precisión de la medición.

**Reflexión y Mejoras**

Me di cuenta de que la selección del instrumento y el cuidado en la medición son cruciales para minimizar la incertidumbre. En futuras mediciones, prestaré más atención a la precisión del instrumento y a la técnica empleada.

---

#### Precisión en la Medida

---

**Comparación de Áreas Medidas**

Hoy comparé la medición del área de un salón y un terreno:
- Salón: (20 ± 1) m², 5% incertidumbre
- Terreno: (200 ± 1) m², 0.5% incertidumbre

Medir el área del salón con valores (5.0 ± 0.1) m y (4.0 ± 0.12) m mostró cómo las incertidumbres porcentuales afectan la precisión del área calculada.

**Medición de Área de Terreno**

Midiendo un terreno con dimensiones:
- Largo: (50.0 ± 0.5) m
- Ancho: (25.0 ± 0.5) m

Calculé el área como 1250 m² con una incertidumbre de 37.5 m². La incertidumbre porcentual fue del 3%. Este ejercicio destacó la importancia de las incertidumbres porcentuales en la precisión de las áreas medidas.

**Medición de Mesas en el Laboratorio**

Medí las mesas del laboratorio:
- Largo promedio: 179.0 cm (±0.2 cm)
- Ancho promedio: 88.6 cm (±0.5 cm)

El área calculada fue de 15,877 cm² con una incertidumbre de 100 cm². Las mediciones repetidas y el cálculo de la desviación estándar me permitieron obtener una mejor comprensión de la incertidumbre asociada a estas medidas.

**Reflexión y Mejoras**

Noté que la precisión del instrumento y la repetición de mediciones son esenciales para reducir la incertidumbre. En futuras mediciones, buscaré instrumentos de mayor precisión y realizaré más repeticiones para mejorar la confiabilidad de los resultados.

---

#### Relación Masa-Volumen y Peso-Masa: Manual vs Automático

---

**Introducción a la Relación Masa-Volumen**

Hoy exploré la relación entre masa y volumen, así como entre peso y masa. La densidad, como propiedad derivada de masa y volumen, se volvió central en este estudio. 

**Experimento con Plastilina y Canicas**

Para medir la masa, utilicé una balanza granataria (±0.01 g), y para medir el volumen, una probeta (±2.5 cm³). Aunque la probeta introdujo una gran incertidumbre, los resultados mostraron una clara relación lineal entre masa y volumen.

**Medición del Peso con un Dinamómetro**

Para medir el peso, usé un dinamómetro (±0.2 N). Las mediciones mostraron una relación lineal clara entre peso y masa, confirmando la proporcionalidad esperada.

**Análisis de Datos y Ajustes Lineales**

Ajusté los datos obtenidos con fórmulas de mejor ajuste para determinar la densidad y la constante de proporcionalidad. Los resultados concordaron bien con los valores teóricos, aunque las discrepancias señalaron posibles errores experimentales.

**Reflexión y Mejoras**

La medición del volumen con una probeta de gran incertidumbre fue un desafío. En el futuro, buscaré métodos de medición de volumen más precisos para reducir la incertidumbre y mejorar la precisión de los resultados.

---

### Conclusión General

Estos experimentos han sido un viaje revelador en la comprensión de la incertidumbre y su impacto en las mediciones. La precisión de los instrumentos y la técnica de medición son cruciales para obtener resultados fiables. La repetición de mediciones y el uso de métodos estadísticos para calcular la incertidumbre son esenciales para interpretar adecuadamente los resultados experimentales.

La experiencia adquirida en estas prácticas me ha proporcionado una base sólida para futuras investigaciones en el campo de la mecánica vectorial y la física experimental. Continuaré refinando mis técnicas y métodos para minimizar la incertidumbre y mejorar la precisión de mis experimentos.

### Diario del Experimento: Movimiento Rectilíneo Uniforme (MRU)

**Introducción al MRU**

Hoy iniciamos un nuevo experimento en el laboratorio de física, enfocándonos en el estudio del movimiento rectilíneo uniforme (MRU). Este tipo de movimiento se caracteriza por una velocidad constante y una trayectoria en línea recta. La tarea principal es investigar la relación proporcional entre la distancia recorrida y el tiempo transcurrido por un deslizador en un riel de aire horizontal. Estoy particularmente interesado en medir la velocidad del deslizador y comparar la precisión de varios métodos de medición.

**Preparativos y Primeros Pasos**

Primero, me encargué de la colocación y nivelación del riel de aire, asegurándome de que todo estuviera perfectamente horizontal con la ayuda de un nivel de albañil. Conecté la compresora de aire al riel y coloqué el deslizador en su lugar. Utilicé una liga para lanzar el deslizador, asegurando que recorriese el riel en aproximadamente dos segundos.

**Mediciones Manuales**

Comencé las mediciones manuales con un cronómetro de mano, registrando los tiempos que el deslizador tardaba en recorrer intervalos de 20 cm. Realicé diez mediciones para cada intervalo y anoté los tiempos cuidadosamente. Me sentí bastante satisfecho con la precisión de mis mediciones manuales, aunque consciente de que el método tiene sus limitaciones en cuanto a exactitud.

**Fotocompuertas y Smart Timer**

A continuación, configuré un par de fotocompuertas, las cuales ofrecían una mayor precisión en la medición de los tiempos. Repetí el procedimiento de lanzamiento y anoté los tiempos obtenidos con el Smart Timer. Esta técnica resultó ser más precisa que el cronómetro manual, ofreciendo datos más consistentes y menos susceptibles a errores humanos.

**Fotografía Estroboscópica**

Luego, intenté utilizar la técnica de fotografía estroboscópica. Configuré el estroboscopio y tomé varias fotos del deslizador en movimiento. Sin embargo, encontré que esta técnica presentaba desafíos significativos. La conversión de revoluciones por minuto a radianes por segundo y la calibración del sistema fotográfico fueron más complicadas de lo esperado. Además, noté algunas inconsistencias en el funcionamiento del estrobo, lo que afectó la precisión de las mediciones.

**Análisis de Resultados**

Después de recopilar los datos, creé tablas y gráficas para analizar la relación entre la distancia y el tiempo. Utilizando el método de regresión lineal, determiné la velocidad del deslizador. Las mediciones manuales con el cronómetro mostraron una velocidad constante de 52.59 cm/s. Las mediciones con fotocompuertas arrojaron una velocidad promedio de 45.34 cm/s, un poco más baja pero aún dentro de un rango aceptable.

La técnica de fotografía estroboscópica, sin embargo, resultó ser la menos confiable debido a los problemas mencionados. Los datos obtenidos de este método mostraron una mayor dispersión y no se alinearon bien con los valores teóricos.

**Reflexión y Mejora**

Al final del día, reflexioné sobre los métodos utilizados y cómo podría mejorar en futuros experimentos. Las mediciones manuales, aunque menos precisas, fueron bastante consistentes. Las fotocompuertas y el Smart Timer demostraron ser herramientas valiosas para obtener datos precisos. En cuanto a la fotografía estroboscópica, reconozco que necesito más práctica y posiblemente mejores equipos para obtener resultados confiables. Quizás un estrobo más moderno y sin fallas sería una inversión útil.

### Diario del Experimento: Movimiento Rectilíneo Uniformemente Acelerado (MRUA)

**Introducción al MRUA**

Hoy nos enfocamos en el estudio del movimiento rectilíneo uniformemente acelerado (MRUA). Este tipo de movimiento se caracteriza por una aceleración constante. En este experimento, inclinamos el riel de aire a distintos ángulos para observar cómo cambia la velocidad del deslizador con el tiempo. Mi objetivo es determinar la relación de proporcionalidad entre la velocidad adquirida y el tiempo transcurrido, además de calcular la aceleración y compararla con los valores teóricos.

**Preparativos y Primeros Pasos**

Primero, nivelé el riel de aire y luego lo incliné a diferentes ángulos (5°, 10°, 15°, 20°). Conecté la compresora de aire y coloqué el deslizador en la parte superior del riel. Utilicé una liga para lanzar el deslizador, asegurando un recorrido uniforme.

**Mediciones y Fotocompuertas**

Comencé midiendo los tiempos de recorrido del deslizador para intervalos de 10 cm utilizando un cronómetro de mano, igual que en el experimento anterior. Luego, configuré las fotocompuertas y el Smart Timer para obtener mediciones más precisas. Noté que, a medida que aumentaba el ángulo de inclinación, el deslizador adquiría mayor velocidad en menor tiempo.

**Fotografía Estroboscópica**

Repetí el procedimiento utilizando la técnica de fotografía estroboscópica. Esta vez, me aseguré de ajustar correctamente las escalas en la aplicación Tracker para minimizar la distorsión del espacio. Sin embargo, el estrobo continuó presentando problemas intermitentes, lo que afectó la consistencia de las fotos.

**Análisis de Resultados**

Al analizar los datos, noté que la aceleración del deslizador aumentaba con el ángulo de inclinación, como se esperaba. Los gráficos de distancia contra tiempo mostraban una pendiente creciente, indicando una aceleración constante. Los datos obtenidos con el cronómetro y las fotocompuertas fueron bastante consistentes y se alinearon bien con los valores teóricos.

En cuanto a la fotografía estroboscópica, los datos seguían siendo inconsistentes. La distorsión del espacio y los problemas con el estrobo afectaron la precisión de las mediciones, haciendo difícil obtener resultados fiables.

**Reflexión y Mejora**

Este experimento reafirmó la importancia de la precisión en las mediciones y la necesidad de herramientas confiables. Las fotocompuertas y el Smart Timer nuevamente demostraron ser muy útiles para obtener datos precisos. La fotografía estroboscópica, aunque interesante, requiere más práctica y posiblemente equipos mejor calibrados. Para futuros experimentos, consideraré invertir en un estrobo más fiable y perfeccionar mi técnica de calibración para minimizar los errores sistemáticos.

En conclusión, ambos experimentos me brindaron una comprensión más profunda de los principios del MRU y MRUA. La experiencia adquirida y las reflexiones sobre los métodos utilizados serán valiosas para futuros estudios en mecánica vectorial. La clave está en la precisión y la consistencia de las mediciones, así como en la continua mejora de las técnicas y equipos utilizados.

### Diario de un Físico: Péndulo Simple y Tiro Parabólico

#### V. Péndulo Simple

**Vivencia del Experimento:**

Hoy me dispuse a entender mejor la relación entre las magnitudes físicas en un péndulo simple. Sabía de antemano que la longitud de la cuerda, el ángulo de lanzamiento (amplitud), el periodo de oscilación, la masa del cuerpo colgante y la gravedad juegan roles importantes. Me intrigaba comprobar si realmente, como la teoría sugiere, la masa del objeto y el ángulo inicial no afectan el periodo de oscilación cuando la longitud de la cuerda es constante.

Con el soporte asegurado, colgué dos péndulos idénticos pero de diferente material, asegurándome de que ambos tuvieran la misma longitud. Al soltarlos desde un ángulo de 5 grados, observé atentamente sus movimientos. Me concentré en detectar cualquier desfase durante las primeras veinte oscilaciones. Repetí el proceso varias veces, ajustando la amplitud a 10, 20 y 30 grados. Era fascinante ver cómo, pese a las diferencias en material y ángulo de lanzamiento, el periodo de oscilación se mantenía constante para pequeñas amplitudes.

Sin embargo, al aumentar la amplitud, noté un desfase progresivo en las oscilaciones. Este efecto anarmónico me llevó a reflexionar sobre cómo las aproximaciones teóricas tienen sus límites en la práctica.

**Reflexión para Mejorar:**

En futuros experimentos, sería ideal utilizar dispositivos de medición más precisos para reducir errores humanos al registrar el tiempo de oscilación. Además, investigar el comportamiento anarmónico a mayores amplitudes podría ofrecer una comprensión más profunda de las limitaciones del modelo ideal de péndulo simple.

---

#### VI. Tiro Parabólico

**Vivencia del Experimento:**

Me embarqué en el estudio del tiro parabólico con el objetivo de determinar la ecuación de movimiento de un proyectil. Equipado con una cámara digital y un software de análisis de movimiento, me preparé para capturar la trayectoria de una pequeña pelota lanzada desde diferentes ángulos.

La preparación fue crucial. Ajusté la sensibilidad y apertura del dispositivo de registro, asegurando que la cámara estuviera perfectamente perpendicular al movimiento del proyectil. Realicé diez disparos para cada uno de los cinco ángulos de lanzamiento, desde 15° hasta 45°, registrando cada trayectoria con precisión.

Al analizar las imágenes, obtuve gráficas de posición contra tiempo que mostraban claramente la trayectoria parabólica del proyectil. Apliqué conceptos de derivada para entender las velocidades y aceleraciones en los ejes x e y. La recta en la gráfica de posición contra tiempo en el eje x indicaba una velocidad constante, mientras que la curva en el eje y reflejaba la aceleración constante debida a la gravedad.

**Reflexión para Mejorar:**

Para mejorar futuros estudios, sería beneficioso considerar la resistencia del aire y otros factores externos que podrían influir en el movimiento del proyectil. Además, perfeccionar el diseño del lanzador para minimizar el recoil y asegurar una mayor precisión en las mediciones ayudaría a obtener resultados aún más fiables.

---

**Conclusión:**

Ambos experimentos me han permitido validar importantes principios de la física clásica. En el caso del péndulo simple, confirmé que el periodo de oscilación depende únicamente de la longitud de la cuerda para pequeñas amplitudes, independientemente de la masa y el ángulo de lanzamiento. En el estudio del tiro parabólico, las predicciones teóricas sobre la trayectoria se vieron mayormente confirmadas, destacando la relación entre la velocidad inicial y el alcance del proyectil.

Estas experiencias me han enseñado la importancia de la precisión en la experimentación y la necesidad de considerar todos los factores posibles para mejorar la validez de los resultados. Continuar perfeccionando las técnicas y reflexionando sobre las discrepancias observadas es esencial para avanzar en la comprensión de los fenómenos físicos.

### Diario del Físico: Estudio del Péndulo Bifilar y Movimiento Circular Uniforme

---

#### Estudio del Péndulo Bifilar

**Experimento: Péndulo Bifilar y Oscilaciones Torsionales**

Hoy, me encuentro inmerso en la fascinante tarea de explorar las oscilaciones torsionales de un péndulo bifilar. El objetivo principal de este estudio es comprender cómo diferentes variables, como la longitud de los hilos, la separación entre ellos y la longitud de la varilla, influyen en el período de oscilación. 

**Montaje del Péndulo Bifilar**

Comencé el experimento armando cuidadosamente el péndulo bifilar. Utilicé una varilla de metal suspendida entre dos soportes sujetos a una base firme. La precisión en la colocación y alineación de los componentes fue crucial para asegurar la validez de las mediciones posteriores.

**Comparación de Períodos**

Inicialmente, comparé los períodos de oscilación de dos péndulos bifilares con varillas de diferentes materiales. Al usar una varilla de madera y otra de metal, observé que la masa de la varilla no influía en el período de oscilación. Este hallazgo refuerza la idea de que, para pequeñas amplitudes, la masa de la varilla es un factor despreciable.

**Influencia de la Amplitud de Oscilación**

Para profundizar en la investigación, examiné cómo la amplitud de oscilación afecta el período. Con amplitudes que variaban entre 5 y 15 grados, noté que las pequeñas variaciones en la amplitud no alteraban significativamente el período de oscilación. Esto confirma la teoría de que, para pequeñas amplitudes, la influencia de la amplitud en el período es mínima.

**Dependencia con la Longitud de los Hilos y la Separación entre Ellos**

Procedí a investigar la dependencia del período con la longitud de los hilos. Al mantener constante la separación de los hilos y variar su longitud, descubrí una relación directa: a mayor longitud de los hilos, mayor es el período de oscilación. 

Posteriormente, analicé cómo la separación entre los hilos afecta el período. Aquí, la relación fue inversa: al aumentar la separación, el período disminuía. Estas observaciones fueron respaldadas por mediciones precisas y repetidas que consolidaron la relación empírica entre estas variables.

**Resultados Cuantitativos**

Los datos recolectados se presentan en tablas y gráficos que muestran claramente las tendencias observadas. Por ejemplo, para una longitud de hilo de 30 cm y una separación de 10 cm, el período registrado fue de 2.2 segundos, mientras que al aumentar la longitud a 40 cm, el período se incrementó a 2.9 segundos.

**Reflexión y Mejora**

Al concluir este conjunto de experimentos, me doy cuenta de la importancia de minimizar los errores de medición y asegurar la precisión en cada paso del procedimiento. En futuras iteraciones, podría mejorar el diseño experimental utilizando dispositivos de medición digital más avanzados para reducir la incertidumbre.

---

#### Movimiento Circular Uniforme (MCU)

**Experimento: Movimiento Circular Uniforme**

Hoy, mi atención se centra en analizar el movimiento circular uniforme de un objeto sobre una mesa de aire. Este experimento tiene como objetivo obtener la ecuación de la trayectoria, la velocidad tangencial y la aceleración centrípeta de un objeto que se mueve en una trayectoria circular.

**Preparación del Experimento**

Comencé nivelando la mesa de aire con un nivel de albañil, asegurándome de que estuviera perfectamente horizontal. Luego, coloqué un disco sobre la mesa y lo fijé a un poste central con un tramo de hilo, minimizando la fricción entre el hilo y el poste. Esta configuración es esencial para asegurar que las fuerzas observadas sean puramente resultado del movimiento circular.

**Medición y Registro de Datos**

Ajusté la longitud del hilo a diferentes valores y conecté la compresora de aire. Al activar la compresora y aplicar un impulso al disco, observé su movimiento circular. Utilicé el software Tracker para registrar y analizar las trayectorias del disco, capturando datos precisos sobre la velocidad tangencial y la aceleración centrípeta.

**Variación del Radio y Masa**

Repetí el experimento modificando el radio de la trayectoria circular en diez valores diferentes. Además, retiré el poste central y pasé un hilo por un orificio, atándolo a un objeto de masa conocida. Esto permitió observar cómo la masa del objeto colgante afecta el movimiento del disco deslizante.

**Resultados y Análisis Estadístico**

Para cada configuración, realicé diez mediciones y promedié los resultados. Los datos recopilados fueron sometidos a un análisis estadístico detallado para determinar las incertidumbres asociadas a cada variable física. Los resultados mostraron que, al aumentar el radio, la velocidad tangencial también aumenta, mientras que la aceleración centrípeta disminuye.

**Reflexión y Mejora**

Al reflexionar sobre este experimento, reconozco que la precisión en la medición de las variables físicas es crucial. Para mejorar futuros experimentos, podría implementar sensores digitales para medir la velocidad y la aceleración de manera más precisa. Además, asegurar que el disco y la mesa estén completamente libres de polvo y otras partículas reducirá aún más la fricción y los errores en las mediciones.

Estos experimentos han sido una experiencia enriquecedora y reveladora. Me han permitido no solo entender mejor los principios de la mecánica vectorial, sino también desarrollar habilidades críticas en el diseño y ejecución de experimentos físicos. Cada reflexión y ajuste en el procedimiento experimental es un paso hacia una mayor precisión y comprensión de los fenómenos físicos.

Fuerzas y Movimiento: Un Estudio Experimental
Introducción y Preparación
Explorar el concepto de fuerza y su relación con el movimiento es una tarea fundamental en la física. Con esta intención, preparé un experimento utilizando un carro de baja fricción y un riel de aire. Después de nivelar el riel y preparar el equipo necesario, comencé con la primera parte del experimento.

Primera Parte: Aplicación de Fuerza Constante
La primera parte del experimento fue desafiante, ya que mantener una fuerza constante sobre el carro resultó complicado. Además, la longitud limitada del riel dificultó las mediciones. Sin embargo, este obstáculo nos permitió aprender valiosas lecciones sobre la precisión y la importancia de contar con el equipo adecuado.

Segunda Parte: Sistema de Dos Masas
En la segunda parte, conecté un hilo de masa despreciable al deslizador, pasando por una polea de baja fricción, con una masa constante al otro extremo. Realicé diez mediciones con diferentes masas en el carro, registrando los datos con precisión. Los resultados fueron congruentes con la teoría, mostrando una relación inversamente proporcional entre la masa del carro y la aceleración experimentada.

Resultados y Reflexión
A pesar de las dificultades iniciales, la segunda parte del experimento fue exitosa. Logré modelar correctamente un sistema de dos masas con una tensión intermediaria. Los datos obtenidos confirmaron nuestras expectativas teóricas, demostrando la relación entre la fuerza aplicada y el movimiento resultante.

Para mejorar, sería esencial contar con un riel de mayor longitud y utilizar dinamómetros más precisos para garantizar la constancia de la fuerza aplicada. Además, la incorporación de técnicas avanzadas de medición podría proporcionar datos aún más exactos.

Estos días en el laboratorio han sido una montaña rusa de desafíos y descubrimientos, recordándome constantemente la belleza y complejidad del mundo físico que nos rodea. Cada experimento, con sus éxitos y dificultades, nos acerca más a una comprensión más profunda de las leyes que rigen nuestro universo.



\end{document}